\documentclass[11pt,a4paper]{article}

% ==================== TEMEL PAKETLER (pdflatex) ====================
\usepackage[T1]{fontenc}
\usepackage[utf8]{inputenc}
\usepackage{lmodern}
\usepackage[a4paper,margin=2cm,top=2.4cm,bottom=2.4cm]{geometry}
\usepackage{xcolor}
\usepackage{colortbl}
\usepackage{array}
\usepackage{booktabs}
\usepackage{longtable}
\usepackage{tabularx}
\usepackage{listings}
\usepackage{fancyhdr}
\usepackage{hyperref}

% ==================== RENK PALETİ ====================
\definecolor{anaMavi}{HTML}{1B3A5B}
\definecolor{vurguMavi}{HTML}{2E6DA4}
\definecolor{acikMavi}{HTML}{EAF2F9}
\definecolor{turuncu}{HTML}{D9701E}
\definecolor{yesil}{HTML}{2E7D32}
\definecolor{acikYesil}{HTML}{E8F5E9}
\definecolor{kirmizi}{HTML}{C0392B}
\definecolor{mor}{HTML}{6C3483}
\definecolor{acikMor}{HTML}{F3E9F7}
\definecolor{acikGri}{HTML}{F4F5F7}
\definecolor{ortaGri}{HTML}{6B7280}
\definecolor{koyuGri}{HTML}{2D2D2D}
\definecolor{kodArka}{HTML}{FBFBFA}

% ==================== HYPERREF ====================
\hypersetup{
  colorlinks=true,
  linkcolor={[HTML]{2E6DA4}},
  urlcolor={[HTML]{2E6DA4}},
  pdftitle={Linux Aygıt Sürücüsü Geliştirme Rehberi},
  pdfauthor={Sürücü Rehberi},
  bookmarksnumbered=true
}

% ==================== BAŞLIK STİLLERİ (titlesec olmadan) ====================
\setcounter{secnumdepth}{0}
\newcounter{secc}
\newcounter{subsecc}[secc]
\newcounter{subsubsecc}[subsecc]
\newcommand{\gsec}[1]{\par\phantomsection\stepcounter{secc}%
  \setcounter{subsecc}{0}%
  \vspace{18pt}{\color{anaMavi}\Large\bfseries\thesecc.\hspace{0.5em}#1}\par
  \vspace{2pt}{\color{turuncu}\rule{\linewidth}{1.6pt}}\par\vspace{8pt}%
  \addcontentsline{toc}{section}{\protect\numberline{\thesecc}#1}\markboth{#1}{#1}}
\newcommand{\gsub}[1]{\par\phantomsection\stepcounter{subsecc}%
  \setcounter{subsubsecc}{0}%
  \vspace{11pt}{\color{vurguMavi}\large\bfseries\thesecc.\thesubsecc\hspace{0.5em}#1}\par\vspace{4pt}%
  \addcontentsline{toc}{subsection}{\protect\numberline{\thesecc.\thesubsecc}#1}}
\newcommand{\gsubsub}[1]{\par\phantomsection\stepcounter{subsubsecc}%
  \vspace{7pt}{\color{koyuGri}\normalsize\bfseries #1}\par\vspace{2pt}}

% Kısım (Part) başlığı
\newcommand{\gpart}[2]{\clearpage
  \begin{center}\vspace*{1.5cm}
  {\color{turuncu}\rule{0.7\linewidth}{1.5pt}}\\[10pt]
  {\color{ortaGri}\Large KISIM #1}\\[8pt]
  {\color{anaMavi}\Huge\bfseries #2}\\[10pt]
  {\color{turuncu}\rule{0.7\linewidth}{1.5pt}}
  \end{center}\vspace{1cm}}

% ==================== BAŞLIK / ALTBİLGİ ====================
\setlength{\headheight}{14pt}
\pagestyle{fancy}
\fancyhf{}
\renewcommand{\headrulewidth}{0.4pt}
\renewcommand{\footrulewidth}{0.4pt}
\fancyhead[L]{\small\color{ortaGri}Linux Aygıt Sürücüsü Geliştirme}
\fancyhead[R]{\small\color{ortaGri}\leftmark}
\fancyfoot[C]{\small\color{ortaGri}\thepage}

% ==================== KOD LİSTELEME (C - çekirdek) ====================
\lstdefinestyle{cstyle}{
  language=C,
  basicstyle=\ttfamily\small\color{koyuGri},
  keywordstyle=\color{vurguMavi}\bfseries,
  commentstyle=\color{yesil}\itshape,
  stringstyle=\color{turuncu},
  numberstyle=\tiny\color{ortaGri},
  numbers=left,
  numbersep=8pt,
  backgroundcolor=\color{kodArka},
  frame=single,
  framesep=6pt,
  rulecolor=\color{ortaGri!40},
  breaklines=true,
  showstringspaces=false,
  tabsize=8,
  columns=fullflexible,
  keepspaces=true,
  extendedchars=true,
  literate={ı}{{\i}}1{İ}{{\.I}}1{ş}{{\c s}}1{Ş}{{\c S}}1{ğ}{{\u g}}1{Ğ}{{\u G}}1{ç}{{\c c}}1{Ç}{{\c C}}1{ö}{{\"o}}1{Ö}{{\"O}}1{ü}{{\"u}}1{Ü}{{\"U}}1,
  morekeywords={size_t,ssize_t,loff_t,dev_t,umode_t,u8,u16,u32,u64,
    s8,s16,s32,s64,uint8_t,uint32_t,__user,__iomem,__init,__exit,
    module_init,module_exit,MODULE_LICENSE,MODULE_AUTHOR,MODULE_DESCRIPTION,
    EXPORT_SYMBOL,bool,true,false,NULL,gfp_t,irqreturn_t,spinlock_t,
    struct,cdev,file_operations,inode,file,device,class,platform_device,
    platform_driver,mutex,atomic_t,wait_queue_head_t,work_struct,
    container_of,likely,unlikely}
}
\lstset{style=cstyle}

% Kabuk/terminal stili
\lstdefinestyle{shstyle}{
  basicstyle=\ttfamily\small\color{koyuGri},
  backgroundcolor=\color{koyuGri!6},
  frame=single,framesep=6pt,rulecolor=\color{ortaGri!40},
  breaklines=true,showstringspaces=false,tabsize=4,columns=fullflexible,
  keepspaces=true,
  literate={ı}{{\i}}1{İ}{{\.I}}1{ş}{{\c s}}1{Ş}{{\c S}}1{ğ}{{\u g}}1{Ğ}{{\u G}}1{ç}{{\c c}}1{Ç}{{\c C}}1{ö}{{\"o}}1{Ö}{{\"O}}1{ü}{{\"u}}1{Ü}{{\"U}}1,
}

% ==================== ÖZEL KUTULAR (tcolorbox olmadan) ====================
\newcommand{\callout}[4]{\par\smallskip\noindent\begingroup
  \setlength{\fboxsep}{8pt}\setlength{\fboxrule}{0.8pt}%
  \fcolorbox{#1}{#2}{\parbox{\dimexpr\linewidth-2\fboxsep-2\fboxrule\relax}{%
    {\bfseries\color{#1}#3}\par\smallskip #4}}\endgroup\par\smallskip}
\newcommand{\bilgi}[1]{\callout{vurguMavi}{acikMavi}{Bilgi}{#1}}
\newcommand{\ipucu}[1]{\callout{yesil}{acikYesil}{İpucu}{#1}}
\newcommand{\uyari}[1]{\callout{kirmizi}{kirmizi!7}{Dikkat}{#1}}
\newcommand{\ornek}[2][Örnek]{\callout{ortaGri!90!black}{acikGri}{#1}{#2}}
\newcommand{\notk}[2][Not]{\callout{mor}{acikMor}{#1}{#2}}

% Yardımcı komutlar
\newcommand{\code}[1]{\texttt{\color{koyuGri}#1}}
\renewcommand{\arraystretch}{1.35}

% ==================== BAŞLANGIÇ ====================
\begin{document}

% ---------- KAPAK ----------
\begin{titlepage}
\centering
\vspace*{1.4cm}
{\color{turuncu}\rule{\textwidth}{2pt}}\\[8pt]
{\color{anaMavi}\Huge\bfseries Linux Aygıt Sürücüsü}\\[10pt]
{\color{vurguMavi}\LARGE Geliştirme Rehberi}\\[8pt]
{\color{ortaGri}\large Çekirdek Modüllerinden Aygıt Modeline Uygulamalı Başvuru}\\[6pt]
{\color{turuncu}\rule{\textwidth}{2pt}}\\[1.4cm]

\begin{center}
\fcolorbox{vurguMavi}{acikMavi}{\parbox{0.86\textwidth}{\centering\large
Bu belge, Linux çekirdek modülü ve aygıt sürücüsü geliştirmeyi
\textbf{çalışan kod örnekleriyle} baştan sona anlatır. İlk ``merhaba dünya''
modülünden karakter aygıtlarına, kesme işleyicilerinden aygıt modeline
ve platform sürücülerine kadar sürücü geliştirmenin temellerini kapsar.\vspace{4pt}}}
\end{center}

\vspace{0.7cm}
\begin{center}
\fcolorbox{ortaGri!50}{acikGri}{\parbox{0.86\textwidth}{\small
\textbf{Kısım 1 --- Çekirdek Modülü Temelleri:} kullanıcı alanı vs çekirdek
\textbullet\ ilk modül \textbullet\ Makefile \& derleme \textbullet\ modül
parametreleri \textbullet\ printk \& günlükleme.\vspace{3pt}

\textbf{Kısım 2 --- Karakter Aygıtları:} aygıt dosyaları \& major/minor
\textbullet\ \code{cdev} \& \code{file\_operations} \textbullet\ open/read/write
\textbullet\ \code{copy\_to/from\_user} \textbullet\ \code{ioctl}
\textbullet\ udev \& otomatik aygıt düğümü.\vspace{3pt}

\textbf{Kısım 3 --- Çekirdek Altyapısı:} bellek ayırma \textbullet\ eşzamanlılık
(spinlock, mutex, atomik) \textbullet\ kesmeler \& alt yarı \textbullet\ zaman
\& zamanlayıcılar.\vspace{3pt}

\textbf{Kısım 4 --- Aygıt Modeli:} sysfs \& kobject \textbullet\ platform
sürücüleri \textbullet\ device tree \textbullet\ kaynak yönetimi.\vspace{3pt}

\textbf{Kısım 5 --- İleri Konular:} bekleme kuyrukları \& bloklamalı G/Ç
\textbullet\ \code{poll} \textbullet\ proc arayüzleri \textbullet\ hata ayıklama.\vspace{3pt}}}
\end{center}

\vfill
{\color{ortaGri}\large Hazırlanma Tarihi: 12 Temmuz 2026}\\[4pt]
{\color{ortaGri} Linux Çekirdek 5.x/6.x \textbullet\ C \textbullet\ Kodlar İngilizce, açıklamalar Türkçe}
\vspace*{0.6cm}
\end{titlepage}

% ---------- İÇİNDEKİLER ----------
\tableofcontents
\thispagestyle{empty}
\clearpage

% #####################################################################
\gpart{1}{Çekirdek Modülü Temelleri}
% #####################################################################

% =====================================================================
\gsec{Giriş: Sürücüler ve Çekirdek Alanı}
% =====================================================================

Aygıt sürücüsü (device driver), işletim sisteminin donanımla konuşmasını
sağlayan yazılım katmanıdır. Uygulamalar donanımı doğrudan bilmez; sürücüler
soyutlama sağlar --- bir uygulama \code{read()} çağırdığında, ister bir
diskten ister bir ağ kartından okusun aynı arayüzü kullanır. Linux'ta
sürücülerin çoğu \emph{çekirdek modülü} olarak yazılır.

\gsub{Kullanıcı Alanı vs Çekirdek Alanı}

Linux belleği iki ayrı bölgeye ayırır. Sürücü yazarken sürekli bu ayrımın
farkında olmalısın, çünkü çekirdekte yapılan bir hata tüm sistemi çökertebilir.

\begin{center}
\begin{tabularx}{\textwidth}{@{}>{\bfseries\color{anaMavi}}l X X@{}}
\toprule
\rowcolor{acikMavi}\color{anaMavi}\textbf{} & \color{anaMavi}\textbf{Kullanıcı Alanı} & \color{anaMavi}\textbf{Çekirdek Alanı}\\
\midrule
Çalışan kod & Uygulamalar & Çekirdek, modüller, sürücüler\\
\rowcolor{acikGri} Ayrıcalık & Kısıtlı (ring 3) & Tam yetki (ring 0)\\
Hata sonucu & Süreç çöker & Sistem çöker (kernel panic)\\
\rowcolor{acikGri} Bellek erişimi & Kendi sanal alanı & Tüm fiziksel bellek\\
Kütüphaneler & libc, tüm kullanıcı kütüph. & Yalnızca çekirdek API'si\\
\bottomrule
\end{tabularx}
\end{center}

\uyari{Çekirdek alanında \textbf{standart C kütüphanesi yoktur}: \code{printf},
\code{malloc}, \code{stdio.h} kullanamazsın. Bunların yerine çekirdek API'si
gelir: \code{printk}, \code{kmalloc}, \code{<linux/*.h>} başlıkları. Ayrıca
kayan nokta (float) işlemleri de kural olarak kullanılmaz. Çekirdek kısıtlı ve
kendine özgü bir ortamdır.}

\gsub{Statik Derleme vs Yüklenebilir Modül}

Bir sürücü çekirdeğe iki şekilde dahil edilebilir:

\begin{center}
\begin{tabularx}{\textwidth}{@{}>{\bfseries\color{anaMavi}}l X@{}}
\toprule
\rowcolor{acikMavi}\color{anaMavi}\textbf{Yöntem} & \color{anaMavi}\textbf{Açıklama}\\
\midrule
Statik (built-in) & Çekirdeğin içine derlenir; her zaman yüklüdür. Yeniden derleme gerekir.\\
\rowcolor{acikGri} Modül (\code{.ko}) & Çalışırken \code{insmod} ile yüklenir, \code{rmmod} ile kaldırılır. Geliştirme için idealdir.\\
\bottomrule
\end{tabularx}
\end{center}

Geliştirme sırasında neredeyse her zaman \emph{yüklenebilir çekirdek modülü}
(Loadable Kernel Module --- LKM) kullanırız; çünkü çekirdeği yeniden
başlatmadan modülü yükleyip kaldırabiliriz.

\gsub{Geliştirme Ortamının Hazırlanması}

Modül derlemek için çalışan çekirdeğinizle eşleşen başlık dosyaları gerekir.

\begin{lstlisting}[style=shstyle]
# Derleme araçları ve çekirdek başlıkları (Debian/Ubuntu)
$ sudo apt install build-essential linux-headers-$(uname -r)

# Fedora/RHEL
$ sudo dnf install kernel-devel kernel-headers gcc make

# Çalışan çekirdek sürümünü gör
$ uname -r                      # ör. 6.8.0-40-generic

# Başlıkların yeri (Makefile bunu kullanır)
$ ls /lib/modules/$(uname -r)/build
\end{lstlisting}

\bilgi{Sürücü geliştirmeyi \textbf{sanal makinede} (QEMU/VirtualBox) yapmak
şiddetle önerilir: çekirdek kodundaki bir hata sistemi çökertebilir veya veri
kaybına yol açabilir. Sanal makinede çökme olsa bile ana sisteminiz güvende
kalır ve anlık görüntülerden (snapshot) hızlıca geri dönebilirsiniz.}

% =====================================================================
\gsec{İlk Çekirdek Modülü}
% =====================================================================

Geleneksel ``merhaba dünya'' ile başlayalım. Bir çekirdek modülünün iki temel
fonksiyonu vardır: yüklendiğinde çalışan bir \emph{giriş} (init) fonksiyonu ve
kaldırıldığında çalışan bir \emph{çıkış} (exit) fonksiyonu.

\gsub{Modülün Kaynak Kodu}

\begin{lstlisting}
#include <linux/module.h>    /* tüm modüller için gerekli */
#include <linux/kernel.h>    /* printk, KERN_INFO */
#include <linux/init.h>      /* __init, __exit makroları */

/* Modül yüklendiğinde çağrılır. __init: bu kod yalnızca bir kez
   çalışır, sonra belleği serbest bırakılabilir. */
static int __init hello_init(void)
{
    printk(KERN_INFO "Merhaba: modül yüklendi\n");
    return 0;                /* 0 = başarı; negatif = hata (yükleme iptal) */
}

/* Modül kaldırıldığında çağrılır. */
static void __exit hello_exit(void)
{
    printk(KERN_INFO "Güle güle: modül kaldırıldı\n");
}

module_init(hello_init);     /* giriş noktasını kaydet */
module_exit(hello_exit);     /* çıkış noktasını kaydet */

MODULE_LICENSE("GPL");       /* lisans; GPL olmayan modüller çekirdeği "kirletir" */
MODULE_AUTHOR("Ad Soyad");
MODULE_DESCRIPTION("Basit bir merhaba dünya modülü");
MODULE_VERSION("1.0");
\end{lstlisting}

\uyari{\code{MODULE\_LICENSE} zorunludur. GPL uyumlu bir lisans belirtmezsen
çekirdek ``tainted'' (kirlenmiş) olarak işaretlenir ve birçok GPL-only çekirdek
sembolüne (fonksiyona) erişemezsin. Ayrıca giriş/çıkış fonksiyonları
\code{static} olmalıdır --- global ad alanını kirletmemek için.}

\gsub{Makefile ve Derleme}

Çekirdek modülleri, çekirdeğin kendi \emph{kbuild} sistemiyle derlenir.
Makefile şaşırtıcı derecede kısadır:

\begin{lstlisting}[style=shstyle]
# Makefile
obj-m += hello.o                 # 'hello.c' -> 'hello.ko' üret

# Çalışan çekirdeğin derleme dizinine yönlendir
KDIR := /lib/modules/$(shell uname -r)/build

all:
	$(MAKE) -C $(KDIR) M=$(PWD) modules

clean:
	$(MAKE) -C $(KDIR) M=$(PWD) clean
\end{lstlisting}

\notk[Kbuild nasıl çalışır?]{\code{obj-m += hello.o} satırı kbuild'e
``\code{hello.c}'yi yüklenebilir bir modül olarak derle'' der. \code{-C \$(KDIR)}
çekirdeğin Makefile'ına geçer, \code{M=\$(PWD)} ise ``modül kaynağım burada''
anlamına gelir. Sonuç \code{.ko} (kernel object) uzantılı bir dosyadır.}

\gsub{Modülü Yükleme ve Test Etme}

\begin{lstlisting}[style=shstyle]
$ make                          # hello.ko üret
$ sudo insmod hello.ko          # modülü yükle -> hello_init çalışır
$ lsmod | grep hello            # yüklü modülleri listele
$ sudo rmmod hello              # modülü kaldır -> hello_exit çalışır

$ dmesg | tail                  # çekirdek günlük mesajlarını gör
# [ ... ] Merhaba: modül yüklendi
# [ ... ] Güle güle: modül kaldırıldı

$ modinfo hello.ko              # modül meta bilgisini gör (lisans, yazar...)
\end{lstlisting}

% =====================================================================
\gsec{Günlükleme ve Modül Parametreleri}
% =====================================================================

\gsub{printk ve Günlük Seviyeleri}

Çekirdekte \code{printf} yerine \code{printk} kullanılır. \code{printk},
mesajlara bir \emph{öncelik seviyesi} ekler; bu seviye mesajın nereye ve
ne zaman gösterileceğini etkiler.

\begin{center}
\begin{tabularx}{\textwidth}{@{}>{\ttfamily\bfseries\color{anaMavi}}l l X@{}}
\toprule
\rowcolor{acikMavi}\color{anaMavi}\textbf{Makro} & \color{anaMavi}\textbf{Seviye} & \color{anaMavi}\textbf{Kullanım}\\
\midrule
KERN\_EMERG & 0 & Sistem kullanılamaz.\\
\rowcolor{acikGri} KERN\_ERR & 3 & Hata durumu.\\
KERN\_WARNING & 4 & Uyarı.\\
\rowcolor{acikGri} KERN\_INFO & 6 & Bilgilendirme.\\
KERN\_DEBUG & 7 & Hata ayıklama.\\
\bottomrule
\end{tabularx}
\end{center}

\begin{lstlisting}
/* Modern çekirdeklerde tercih edilen yardımcılar (pr_* makroları) */
pr_info("aygıt %d başlatıldı\n", id);       /* = printk(KERN_INFO ...) */
pr_err("bellek ayrılamadı\n");              /* = printk(KERN_ERR ...) */
pr_warn("tampon neredeyse dolu\n");
pr_debug("hata ayıklama: sayaç=%d\n", count); /* yalnızca DEBUG etkinse */

/* Bir sürücü içinde daha iyisi: dev_* (aygıtı da yazdırır) */
dev_info(dev, "başlatma tamam\n");
dev_err(dev, "kayıt başarısız: %d\n", ret);
\end{lstlisting}

\ipucu{\code{dmesg} çekirdek halka tamponunu (ring buffer) gösterir; sürücü
geliştirmede en çok kullanacağın araçtır. Canlı izlemek için \code{dmesg -w}
(watch) kullan. \code{printk} biçim belirteçleri \code{printf}'e benzer ama
çekirdeğe özgü olanlar da vardır: işaretçi için \code{\%p}, güvenli hash'li
işaretçi; fiziksel adres için \code{\%pa}; \code{resource} için \code{\%pr}.}

\gsub{Modül Parametreleri}

Modüllere, yükleme anında dışarıdan değer geçirebiliriz --- tıpkı bir
programa komut satırı argümanı geçmek gibi. Bu, aynı modülü farklı ayarlarla
kullanmayı sağlar.

\begin{lstlisting}
#include <linux/moduleparam.h>

static int buffer_size = 1024;              /* varsayılan değer */
static char *device_name = "mydev";
static int debug_enabled = 0;

/* Parametreyi kaydet: (değişken, tip, sysfs izinleri) */
module_param(buffer_size, int, 0644);
MODULE_PARM_DESC(buffer_size, "Tampon boyutu (bayt)");

module_param(device_name, charp, 0644);     /* charp = char pointer */
MODULE_PARM_DESC(device_name, "Aygıt adı");

module_param(debug_enabled, bool, 0644);
MODULE_PARM_DESC(debug_enabled, "Hata ayıklama günlüğünü aç (0/1)");

static int __init my_init(void)
{
    pr_info("tampon=%d ad=%s hata_ayıkla=%d\n",
            buffer_size, device_name, debug_enabled);
    return 0;
}
\end{lstlisting}

\begin{lstlisting}[style=shstyle]
# Parametreleri yükleme anında geç
$ sudo insmod mymod.ko buffer_size=4096 device_name="sensor" debug_enabled=1

# Yüklü bir modülün parametrelerini sysfs'ten oku (0644 izinleri sayesinde)
$ cat /sys/module/mymod/parameters/buffer_size
4096
\end{lstlisting}

\bilgi{\code{module\_param}'ın üçüncü argümanı, parametrenin
\code{/sys/module/.../parameters/} altındaki dosya izinleridir. \code{0644}
ile parametre çalışırken okunabilir (ve \code{root} tarafından yazılabilir);
\code{0} verirsen sysfs'te hiç görünmez. Böylece bir sürücünün davranışını
yeniden yüklemeden ayarlayabilirsin.}

% #####################################################################
\gpart{2}{Karakter Aygıtları}
% #####################################################################

% =====================================================================
\gsec{Aygıt Dosyaları ve Aygıt Numaraları}
% =====================================================================

Linux'ta ``her şey bir dosyadır'' ilkesi donanım için de geçerlidir.
Aygıtlar, \code{/dev} altındaki özel dosyalar aracılığıyla kullanıcı alanına
sunulur. Bir uygulama bu dosyaya \code{open}, \code{read}, \code{write}
yaptığında, çekirdek çağrıyı ilgili sürücüye yönlendirir.

\gsub{Karakter vs Blok Aygıtları}

\begin{center}
\begin{tabularx}{\textwidth}{@{}>{\bfseries\color{anaMavi}}l X X@{}}
\toprule
\rowcolor{acikMavi}\color{anaMavi}\textbf{Tür} & \color{anaMavi}\textbf{Erişim} & \color{anaMavi}\textbf{Örnek}\\
\midrule
Karakter (char) & Bayt akışı, sırayla & Seri port, klavye, \code{/dev/null}\\
\rowcolor{acikGri} Blok (block) & Sabit boyutlu bloklar, rastgele & Disk, SSD, USB bellek\\
Ağ (network) & Paket, soket arayüzü & Ethernet, Wi-Fi\\
\bottomrule
\end{tabularx}
\end{center}

Bu rehber \emph{karakter aygıtlarına} odaklanır; çünkü en yaygın ve öğrenmesi
en kolay sürücü türüdür. Sensörler, GPIO, seri portlar genelde karakter
aygıtı olarak sunulur.

\gsub{Major ve Minor Numaraları}

Her aygıt dosyası iki sayıyla tanımlanır: \emph{major} numara hangi sürücünün
sorumlu olduğunu, \emph{minor} numara ise o sürücünün hangi somut aygıtı
(birden çok aynı tip aygıt olabilir) belirtir.

\begin{lstlisting}[style=shstyle]
$ ls -l /dev/null /dev/ttyS0
crw-rw-rw- 1 root root 1, 3 ... /dev/null    # major=1, minor=3
crw-rw---- 1 root dialout 4, 64 ... /dev/ttyS0  # major=4, minor=64
#  ^ 'c' = karakter aygıtı (block için 'b')
\end{lstlisting}

\begin{lstlisting}
#include <linux/fs.h>

dev_t dev;                                  /* aygıt numarası (major+minor) */

/* Çekirdekten dinamik olarak bir major numara iste (önerilen) */
ret = alloc_chrdev_region(&dev, 0, 1, "mydev");
/*   &dev: sonuç, 0: ilk minor, 1: aygıt sayısı, "mydev": ad */

int major = MAJOR(dev);                     /* numaradan major çıkar */
int minor = MINOR(dev);                     /* numaradan minor çıkar */
pr_info("major=%d minor=%d\n", major, minor);

/* Kullanım sonrası serbest bırak (exit fonksiyonunda) */
unregister_chrdev_region(dev, 1);
\end{lstlisting}

\ipucu{Eskiden sabit major numaralar (\code{register\_chrdev}) kullanılırdı ama
çakışma riski vardır. Modern sürücüler \code{alloc\_chrdev\_region} ile
çekirdekten \emph{dinamik} numara ister --- böylece sistemdeki başka bir
sürücüyle çakışmaz.}

% =====================================================================
\gsec{Karakter Sürücüsü: file\_operations}
% =====================================================================

Bir karakter sürücüsünün kalbi \code{struct file\_operations}'tır: kullanıcı
alanı sistem çağrılarını (\code{open}, \code{read}, \code{write}\dots) senin
sürücü fonksiyonlarına bağlayan bir işaretçi tablosudur. Çekirdek, uygulama
aygıt dosyasına bir işlem yaptığında ilgili fonksiyonu çağırır.

\gsub{file\_operations Yapısı}

\begin{lstlisting}
#include <linux/fs.h>

static int     my_open(struct inode *inode, struct file *file);
static int     my_release(struct inode *inode, struct file *file);
static ssize_t my_read(struct file *file, char __user *buf,
                       size_t count, loff_t *offset);
static ssize_t my_write(struct file *file, const char __user *buf,
                        size_t count, loff_t *offset);

static const struct file_operations my_fops = {
    .owner   = THIS_MODULE,     /* modül sayacı yönetimi için */
    .open    = my_open,
    .release = my_release,      /* close() sistem çağrısına karşılık gelir */
    .read    = my_read,
    .write   = my_write,
    .llseek  = default_llseek,
};
\end{lstlisting}

\notk[\code{\_\_user} ne anlama gelir?]{\code{char \_\_user *buf} imzasındaki
\code{\_\_user}, bu işaretçinin \emph{kullanıcı alanına} ait olduğunu belirtir.
Çekirdek bu işaretçiyi doğrudan dereference \textbf{edemez} --- geçerliliği
garanti değildir ve güvenlik açığı olur. Bunun yerine \code{copy\_to\_user} /
\code{copy\_from\_user} kullanılır. \code{sparse} statik analiz aracı bu
işaretlemeyi denetler.}

\gsub{cdev ile Aygıtı Kaydetme}

Aygıt numarasını aldıktan sonra, \code{struct cdev} ile \code{file\_operations}
tablosunu çekirdeğe kaydederiz. Böylece o major/minor'a yapılan çağrılar
bizim fonksiyonlarımıza gelir.

\begin{lstlisting}
#include <linux/cdev.h>

static struct cdev my_cdev;                 /* karakter aygıtı yapısı */

static int __init my_init(void)
{
    int ret;

    /* 1) Aygıt numarası ayır */
    ret = alloc_chrdev_region(&dev, 0, 1, "mydev");
    if (ret < 0) {
        pr_err("aygıt numarası ayrılamadı\n");
        return ret;
    }

    /* 2) cdev'i file_operations ile ilklendir ve kaydet */
    cdev_init(&my_cdev, &my_fops);
    my_cdev.owner = THIS_MODULE;
    ret = cdev_add(&my_cdev, dev, 1);       /* çekirdeğe ekle */
    if (ret < 0) {
        unregister_chrdev_region(dev, 1);   /* hata: geri al */
        return ret;
    }

    pr_info("sürücü hazır (major=%d)\n", MAJOR(dev));
    return 0;
}

static void __exit my_exit(void)
{
    cdev_del(&my_cdev);                     /* kaydı kaldır */
    unregister_chrdev_region(dev, 1);       /* numarayı serbest bırak */
}
\end{lstlisting}

\uyari{Çekirdek programlamada \textbf{hata yönetimi ve temizlik} kritiktir.
Bir adım başarısız olursa, o ana kadar ayrılan tüm kaynakları
\emph{ters sırayla} serbest bırakmalısın (yukarıdaki örnekte
\code{alloc\_chrdev\_region} başarılı ama \code{cdev\_add} başarısızsa,
numarayı geri veriyoruz). Sızdırılan çekirdek kaynağı, yeniden başlatmaya
kadar geri gelmez.}

% =====================================================================
\gsec{Okuma, Yazma ve Kullanıcı Verisi}
% =====================================================================

Sürücünün asıl işi \code{read} ve \code{write} fonksiyonlarındadır. Buradaki
en önemli kural: çekirdek ile kullanıcı alanı arasında veri, asla doğrudan
işaretçiyle değil, özel fonksiyonlarla taşınır.

\gsub{copy\_to\_user ve copy\_from\_user}

\begin{lstlisting}
#include <linux/uaccess.h>      /* copy_to_user, copy_from_user */
#include <linux/slab.h>         /* kmalloc, kfree */

#define BUF_SIZE 1024
static char *device_buffer;     /* çekirdek içi tampon */
static size_t buffer_len;       /* içindeki geçerli bayt sayısı */

/* Kullanıcı aygıttan OKUYOR: veri çekirdek -> kullanıcı yönünde akar */
static ssize_t my_read(struct file *file, char __user *buf,
                       size_t count, loff_t *offset)
{
    size_t remaining;

    if (*offset >= buffer_len)              /* dosya sonu */
        return 0;

    /* İstenen kadar ama en fazla kalan veri kadar oku */
    if (count > buffer_len - *offset)
        count = buffer_len - *offset;

    /* Çekirdek tamponundan kullanıcı tamponuna kopyala.
       Dönüş: KOPYALANAMAYAN bayt sayısı (0 = tam başarı). */
    remaining = copy_to_user(buf, device_buffer + *offset, count);
    if (remaining)
        return -EFAULT;                     /* geçersiz kullanıcı adresi */

    *offset += count;                       /* konumu ilerlet */
    pr_debug("%zu bayt okundu\n", count);
    return count;                           /* okunan bayt sayısını döndür */
}

/* Kullanıcı aygıta YAZIYOR: veri kullanıcı -> çekirdek yönünde akar */
static ssize_t my_write(struct file *file, const char __user *buf,
                        size_t count, loff_t *offset)
{
    if (count > BUF_SIZE)                    /* taşmayı önle */
        count = BUF_SIZE;

    if (copy_from_user(device_buffer, buf, count))
        return -EFAULT;

    buffer_len = count;
    pr_debug("%zu bayt yazıldı\n", count);
    return count;
}
\end{lstlisting}

\uyari{\code{copy\_to\_user}/\code{copy\_from\_user} \textbf{başarısızlıkta
kopyalanamayan bayt sayısını} döndürür, hata kodu değil. Sıfır dönmesi tam
başarı demektir. Ayrıca bu fonksiyonlar \emph{uyuyabilir} (sayfa hatası
tetikleyebilir), bu yüzden atomik bağlamda (spinlock tutarken, kesme
işleyicide) çağrılamazlar. Doğrudan \code{*buf = x} yapmak ise ciddi bir
güvenlik açığı ve çökme sebebidir.}

\gsub{open ve release}

\begin{lstlisting}
static int my_open(struct inode *inode, struct file *file)
{
    /* Tampon henüz yoksa ayır (basit örnek; genelde init'te yapılır) */
    if (!device_buffer) {
        device_buffer = kmalloc(BUF_SIZE, GFP_KERNEL);
        if (!device_buffer)
            return -ENOMEM;                 /* bellek yok */
    }
    pr_info("aygıt açıldı\n");
    return 0;
}

static int my_release(struct inode *inode, struct file *file)
{
    pr_info("aygıt kapatıldı\n");
    return 0;
}
\end{lstlisting}

\gsub{Sürücüyü Test Etme}

\begin{lstlisting}[style=shstyle]
# Aygıt düğümünü elle oluştur (major numarayı dmesg'den al)
$ sudo mknod /dev/mydev c 240 0     # c=karakter, 240=major, 0=minor
$ sudo chmod 666 /dev/mydev

# Yaz ve oku
$ echo "merhaba sürücü" > /dev/mydev
$ cat /dev/mydev
merhaba sürücü

# Düşük seviye test
$ dd if=/dev/mydev bs=1 count=10
\end{lstlisting}

% =====================================================================
\gsec{ioctl: Aygıt Kontrol Komutları}
% =====================================================================

\code{read}/\code{write} bayt akışı içindir, ama aygıtlara çoğu zaman
\emph{komut} göndermemiz gerekir: ``baud hızını ayarla'', ``sıfırla'',
``durumu sorgula''. Bunun için \code{ioctl} (input/output control) kullanılır.

\gsub{ioctl Komutlarını Tanımlama}

Komut numaraları rastgele seçilmez; çakışmayı önleyen makrolarla üretilir.
Bu makrolar komuta bir yön (okuma/yazma) ve veri boyutu da kodlar.

\begin{lstlisting}
#include <linux/ioctl.h>

#define MYDEV_MAGIC 'k'                     /* sürücüye özel "sihirli" harf */

/* _IO: veri yok, _IOR: oku, _IOW: yaz, _IOWR: her ikisi */
#define MYDEV_RESET     _IO(MYDEV_MAGIC, 0)
#define MYDEV_GET_SIZE  _IOR(MYDEV_MAGIC, 1, int)
#define MYDEV_SET_SIZE  _IOW(MYDEV_MAGIC, 2, int)

static long my_ioctl(struct file *file, unsigned int cmd, unsigned long arg)
{
    int value;

    switch (cmd) {
    case MYDEV_RESET:
        buffer_len = 0;                     /* tamponu sıfırla */
        memset(device_buffer, 0, BUF_SIZE);
        pr_info("aygıt sıfırlandı\n");
        break;

    case MYDEV_GET_SIZE:                    /* çekirdek -> kullanıcı */
        value = buffer_len;
        if (copy_to_user((int __user *)arg, &value, sizeof(value)))
            return -EFAULT;
        break;

    case MYDEV_SET_SIZE:                    /* kullanıcı -> çekirdek */
        if (copy_from_user(&value, (int __user *)arg, sizeof(value)))
            return -EFAULT;
        if (value < 0 || value > BUF_SIZE)
            return -EINVAL;                 /* geçersiz argüman */
        buffer_len = value;
        break;

    default:
        return -ENOTTY;                     /* bilinmeyen komut */
    }
    return 0;
}

/* file_operations'a ekle: */
/*   .unlocked_ioctl = my_ioctl, */
\end{lstlisting}

\gsub{Kullanıcı Alanından ioctl Çağırma}

\begin{lstlisting}
/* Kullanıcı alanı programı (aynı komut tanımlarını paylaşır) */
#include <sys/ioctl.h>
#include <fcntl.h>

int fd = open("/dev/mydev", O_RDWR);

ioctl(fd, MYDEV_RESET);                     /* argümansız komut */

int size = 512;
ioctl(fd, MYDEV_SET_SIZE, &size);           /* değer gönder */

int current;
ioctl(fd, MYDEV_GET_SIZE, &current);        /* değer al */
printf("mevcut boyut: %d\n", current);

close(fd);
\end{lstlisting}

\bilgi{\code{ioctl} güçlüdür ama tip güvenliği zayıftır (her şey
\code{unsigned long arg}'dan geçer). Modern çekirdek geliştirmede, basit
durum sorguları için \code{ioctl} yerine \code{sysfs} nitelikleri (Kısım 4)
tercih edilir --- çünkü kabuktan \code{cat}/\code{echo} ile erişilebilir ve
kendini belgeler. Yine de karmaşık, ikili protokoller için \code{ioctl}
vazgeçilmezdir.}

% =====================================================================
\gsec{Otomatik Aygıt Düğümü: udev ve class}
% =====================================================================

Yukarıda aygıt düğümünü elle \code{mknod} ile oluşturduk --- bu zahmetli ve
hataya açıktır. Gerçek sürücüler, çekirdeğin \emph{aygıt modeli} üzerinden bir
\code{class} ve \code{device} kaydeder; \emph{udev} arka planı bunu görüp
\code{/dev} altında düğümü \textbf{otomatik} oluşturur.

\begin{lstlisting}
#include <linux/device.h>

static struct class *my_class;
static struct device *my_device;

static int __init my_init(void)
{
    /* ... alloc_chrdev_region + cdev_add (yukarıdaki gibi) ... */

    /* 1) Bir aygıt sınıfı oluştur -> /sys/class/myclass/ */
    my_class = class_create("myclass");     /* çekirdek 6.4+ imzası */
    if (IS_ERR(my_class)) {
        cdev_del(&my_cdev);
        unregister_chrdev_region(dev, 1);
        return PTR_ERR(my_class);
    }

    /* 2) Aygıtı oluştur -> udev /dev/mydev düğümünü otomatik yapar */
    my_device = device_create(my_class, NULL, dev, NULL, "mydev");
    if (IS_ERR(my_device)) {
        class_destroy(my_class);
        cdev_del(&my_cdev);
        unregister_chrdev_region(dev, 1);
        return PTR_ERR(my_device);
    }

    pr_info("/dev/mydev otomatik oluşturuldu\n");
    return 0;
}

static void __exit my_exit(void)
{
    device_destroy(my_class, dev);          /* ters sırayla temizle */
    class_destroy(my_class);
    cdev_del(&my_cdev);
    unregister_chrdev_region(dev, 1);
}
\end{lstlisting}

\notk[IS\_ERR ve PTR\_ERR]{Birçok çekirdek fonksiyonu, hata durumunda
\code{NULL} yerine bir \emph{hata işaretçisi} döndürür (adres uzayının en
üstündeki özel değerler). \code{IS\_ERR(ptr)} bunun bir hata olup olmadığını
kontrol eder, \code{PTR\_ERR(ptr)} ise onu bir hata koduna (\code{-ENOMEM}
gibi) çevirir. Bu, çekirdekte çok yaygın bir kalıptır; işaretçi dönen
fonksiyonlarda mutlaka kontrol et.}

\bilgi{Bu noktada tam işlevsel bir karakter sürücün var: dinamik major numara,
\code{cdev} kaydı, \code{read}/\code{write}/\code{ioctl} ve \code{udev} ile
otomatik \code{/dev} düğümü. Bu iskelet, gerçek sürücülerin çoğunun temelini
oluşturur --- geri kalanı, \code{read}/\code{write} içinde gerçek donanımla
konuşmaktan ibarettir.}

% #####################################################################
\gpart{3}{Çekirdek Altyapısı}
% #####################################################################

% =====================================================================
\gsec{Çekirdekte Bellek Yönetimi}
% =====================================================================

Çekirdekte \code{malloc} yoktur; bellek ayırmanın kendine özgü fonksiyonları
ve kuralları vardır. Yanlış ayırma bayrağı veya sızıntı, tüm sistemi
etkileyebilir --- kullanıcı alanındaki gibi süreç bitince otomatik temizlik
yoktur.

\gsub{kmalloc, kzalloc ve vmalloc}

\begin{center}
\begin{tabularx}{\textwidth}{@{}>{\ttfamily\bfseries\color{anaMavi}}l X@{}}
\toprule
\rowcolor{acikMavi}\color{anaMavi}\textbf{Fonksiyon} & \color{anaMavi}\textbf{Özellik}\\
\midrule
kmalloc & Fiziksel olarak \emph{bitişik} bellek; küçük, hızlı; DMA'ya uygun.\\
\rowcolor{acikGri} kzalloc & \code{kmalloc} + sıfırlama (önerilir).\\
vmalloc & Sanal olarak bitişik (fiziksel dağınık); büyük ayırmalar için.\\
\rowcolor{acikGri} kfree / vfree & İlgili serbest bırakma fonksiyonları.\\
\bottomrule
\end{tabularx}
\end{center}

\begin{lstlisting}
#include <linux/slab.h>

/* GFP_KERNEL: normal ayırma; uyuyabilir (bloklayabilir).
   GFP_ATOMIC: uyumaz; kesme bağlamında / spinlock tutarken kullan. */
char *buffer = kmalloc(size, GFP_KERNEL);
if (!buffer)
    return -ENOMEM;                         /* ayırma başarısız */

/* Sıfırlanmış bellek (calloc benzeri) -- tercih edilen */
struct my_data *data = kzalloc(sizeof(*data), GFP_KERNEL);
if (!data)
    return -ENOMEM;

/* Büyük tampon (fiziksel bitişiklik gerekmiyorsa) */
void *big = vmalloc(10 * 1024 * 1024);      /* 10 MB */

/* Serbest bırak (doğru eşleşen fonksiyonla!) */
kfree(buffer);
kfree(data);
vfree(big);
\end{lstlisting}

\uyari{\textbf{GFP bayrağını doğru seç:} \code{GFP\_KERNEL} bellek beklerken
\emph{uyuyabilir}; bu yüzden kesme işleyicide veya spinlock tutarken
kullanılamaz --- orada \code{GFP\_ATOMIC} gerekir (uyumaz ama başarısız olma
olasılığı daha yüksektir). Yanlış bağlamda \code{GFP\_KERNEL} kullanmak
kilitlenmeye (deadlock) yol açar.}

\gsub{Kaynak Yönetimli (devm) Ayırma}

Modern sürücüler, hata yollarını basitleştirmek için \code{devm\_*}
fonksiyonlarını kullanır: bunlar aygıta ``bağlı'' ayırmalardır ve aygıt
kaldırıldığında \emph{otomatik} serbest bırakılır. Böylece \code{kfree}
çağrılarını unutma riski ortadan kalkar.

\begin{lstlisting}
/* Aygıt kaldırıldığında otomatik serbest bırakılır -- kfree gerekmez! */
struct my_data *data = devm_kzalloc(dev, sizeof(*data), GFP_KERNEL);
if (!data)
    return -ENOMEM;
/* ... temizlik kodunda kfree(data) YOK; çekirdek halleder */
\end{lstlisting}

% =====================================================================
\gsec{Eşzamanlılık ve Kilitleme}
% =====================================================================

Çekirdek doğası gereği eşzamanlıdır: aynı sürücü kodu birden çok CPU'da,
birden çok süreç tarafından, hatta bir kesmeyle yarıda kesilerek aynı anda
çalışabilir. Paylaşılan veriyi korumadan bırakırsan \emph{yarış koşulları}
(race conditions) oluşur --- veri bozulur, sistem çöker. Kilitleme bunu önler.

\gsub{Yarış Koşulu Neden Oluşur?}

\code{counter++} gibi ``basit'' bir işlem bile aslında üç adımdır (oku, artır,
yaz). İki CPU aynı anda yaparsa güncellemelerden biri kaybolabilir. Çözüm:
kritik bölgeyi (paylaşılan veriye eriştiğin kod) bir kilitle korumak.

\begin{center}
\begin{tabularx}{\textwidth}{@{}>{\ttfamily\bfseries\color{anaMavi}}l X@{}}
\toprule
\rowcolor{acikMavi}\color{anaMavi}\textbf{Mekanizma} & \color{anaMavi}\textbf{Ne zaman?}\\
\midrule
mutex & Uyuyabilen bağlam (süreç); kritik bölge uzun/uyuyabilir.\\
\rowcolor{acikGri} spinlock & Kesme bağlamı veya kısa kritik bölge; uyuyamaz.\\
atomic\_t & Tek bir tam sayı sayaç/bayrak; kilitsiz.\\
\rowcolor{acikGri} RCU & Çok okuyan, az yazan; yüksek performans.\\
\bottomrule
\end{tabularx}
\end{center}

\gsub{Mutex}

Mutex, kritik bölgeyi tutan iş parçacığı uyuyabildiğinde kullanılır (örneğin
içinde \code{copy\_to\_user} veya \code{kmalloc(GFP\_KERNEL)} olan bölge).

\begin{lstlisting}
#include <linux/mutex.h>

static DEFINE_MUTEX(my_lock);               /* statik mutex tanımla ve ilklendir */

static ssize_t my_write(struct file *file, const char __user *buf,
                        size_t count, loff_t *offset)
{
    if (mutex_lock_interruptible(&my_lock))  /* kilidi al (sinyalle kesilebilir) */
        return -ERESTARTSYS;

    /* --- kritik bölge: paylaşılan tampona güvenle eriş --- */
    if (copy_from_user(device_buffer, buf, count)) {
        mutex_unlock(&my_lock);             /* hata yolunda da kilidi bırak! */
        return -EFAULT;
    }
    buffer_len = count;
    /* --- kritik bölge sonu --- */

    mutex_unlock(&my_lock);                 /* kilidi bırak */
    return count;
}
\end{lstlisting}

\gsub{Spinlock}

Spinlock, kilidi alamayan CPU'nun \emph{uyumadan bekleyip döndüğü} (busy-wait)
bir kilittir. Kesme işleyicilerinde veya çok kısa kritik bölgelerde kullanılır
--- çünkü orada uyumak yasaktır.

\begin{lstlisting}
#include <linux/spinlock.h>

static DEFINE_SPINLOCK(my_spinlock);
static int shared_counter;

void update_counter(void)
{
    unsigned long flags;

    /* irqsave: kilidi al VE bu CPU'daki kesmeleri kapat.
       Kesmeyle paylaşılan veri için şart (deadlock önler). */
    spin_lock_irqsave(&my_spinlock, flags);

    shared_counter++;                       /* çok kısa kritik bölge */

    spin_unlock_irqrestore(&my_spinlock, flags);  /* bırak + kesmeleri geri aç */
}
\end{lstlisting}

\uyari{\textbf{Spinlock tutarken asla uyuma!} Kilit tutulurken
\code{kmalloc(GFP\_KERNEL)}, \code{copy\_to\_user}, \code{mutex\_lock} gibi
uyuyabilen hiçbir çağrı yapılamaz --- yaparsan kilitlenme olur. Ayrıca kritik
bölgeyi olabildiğince kısa tut, çünkü bekleyen CPU boşa döner (güç ve
performans israfı).}

\gsub{Atomik İşlemler}

Sadece tek bir sayacı korumak için tam bir kilit gereksiz yüktür. Atomik
tipler, bölünmez okuma-değiştirme-yazma işlemleri sunar.

\begin{lstlisting}
#include <linux/atomic.h>

static atomic_t open_count = ATOMIC_INIT(0);

static int my_open(struct inode *inode, struct file *file)
{
    atomic_inc(&open_count);                /* kilitsiz, güvenli artırma */
    pr_info("aygıt %d kez açık\n", atomic_read(&open_count));
    return 0;
}

static int my_release(struct inode *inode, struct file *file)
{
    atomic_dec(&open_count);
    return 0;
}
\end{lstlisting}

% =====================================================================
\gsec{Kesmeler (Interrupts)}
% =====================================================================

Donanım, ilgi istediğinde (veri hazır, tuşa basıldı\dots) CPU'ya bir
\emph{kesme} (interrupt) sinyali gönderir. Sürücü, bu kesmeye bir
\emph{işleyici} (handler --- ISR) kaydeder. Kesme geldiğinde çekirdek, o an ne
yapıyorsa duraklatıp işleyiciyi çalıştırır. Bu, donanımı sürekli yoklamaktan
(polling) çok daha verimlidir.

\gsub{Kesme İşleyicisi Kaydetme}

\begin{lstlisting}
#include <linux/interrupt.h>

static int irq_number;

/* Kesme işleyicisi (ISR). ATOMIK bağlamda çalışır: uyuyamaz,
   kmalloc(GFP_KERNEL)/copy_to_user çağıramaz, kısa olmalı. */
static irqreturn_t my_irq_handler(int irq, void *dev_id)
{
    /* En kısa işi burada yap (donanımı onayla, bayrak set et) */
    pr_debug("kesme alındı: irq=%d\n", irq);

    return IRQ_HANDLED;                     /* kesme bizim; işlendi */
}

static int __init my_init(void)
{
    irq_number = 42;                        /* gerçekte donanımdan/DT'den alınır */

    /* İşleyiciyi kaydet */
    int ret = request_irq(irq_number, my_irq_handler,
                          IRQF_SHARED,      /* IRQ hattı paylaşılabilir */
                          "mydev", &my_dev_id);
    if (ret)
        return ret;
    return 0;
}

static void __exit my_exit(void)
{
    free_irq(irq_number, &my_dev_id);       /* işleyiciyi kaldır */
}
\end{lstlisting}

\gsub{Üst Yarı ve Alt Yarı (Top/Bottom Half)}

Kesme işleyicisi kısa olmalıdır --- çalışırken diğer kesmeler gecikir. Bu
yüzden iş ikiye bölünür: \textbf{üst yarı} (top half --- gerçek ISR) yalnızca
acil, kısa işi yapar; ağır iş \textbf{alt yarıya} (bottom half) ertelenir ve
daha güvenli bir bağlamda çalışır.

\begin{center}
\begin{tabularx}{\textwidth}{@{}>{\bfseries\color{anaMavi}}l X X@{}}
\toprule
\rowcolor{acikMavi}\color{anaMavi}\textbf{Mekanizma} & \color{anaMavi}\textbf{Bağlam} & \color{anaMavi}\textbf{Uyuyabilir mi?}\\
\midrule
Tasklet & Atomik (softirq) & Hayır (eski, artık workqueue tercih edilir)\\
\rowcolor{acikGri} Workqueue & Süreç (kernel thread) & Evet\\
Threaded IRQ & Kendi çekirdek ipliği & Evet\\
\bottomrule
\end{tabularx}
\end{center}

\begin{lstlisting}
#include <linux/workqueue.h>

/* Alt yarı: süreç bağlamında çalışır -> uyuyabilir, kmalloc yapabilir */
static void my_work_handler(struct work_struct *work)
{
    /* Ağır işi burada güvenle yap (uzun işlem, uyuyan çağrılar...) */
    pr_info("alt yarı çalışıyor: veri işleniyor\n");
}

static DECLARE_WORK(my_work, my_work_handler);

/* Üst yarı (ISR): sadece işi zamanla, hemen dön */
static irqreturn_t my_irq_handler(int irq, void *dev_id)
{
    schedule_work(&my_work);                /* alt yarıyı tetikle */
    return IRQ_HANDLED;
}
\end{lstlisting}

\bilgi{\textbf{Threaded IRQ} modern ve pratik bir yaklaşımdır:
\code{request\_threaded\_irq()} ile hem hızlı bir üst yarı hem de kendi
ipliğinde uyuyabilen bir alt yarı fonksiyonu verirsin. Çekirdek ikisini
otomatik yönetir; ayrı workqueue/tasklet kurmana gerek kalmaz.}

% =====================================================================
\gsec{Çekirdekte Zaman}
% =====================================================================

Sürücüler sık sık zamanla ilgilenir: gecikme koyma, zaman aşımı, periyodik iş.
Çekirdeğin zaman birimi \emph{jiffy}'dir --- sistem açılışından beri geçen
zamanlayıcı tik sayısı.

\gsub{jiffies ve Gecikmeler}

\begin{lstlisting}
#include <linux/jiffies.h>
#include <linux/delay.h>

unsigned long start = jiffies;              /* şu anki tik sayısı */

/* HZ = saniyedeki tik sayısı (genelde 100, 250 veya 1000) */
unsigned long timeout = jiffies + 2 * HZ;   /* şimdiden 2 saniye sonrası */

if (time_after(jiffies, timeout))           /* zaman aşımı kontrolü */
    pr_warn("zaman aşımı!\n");

/* Kısa MEŞGUL bekleme (uyumaz) -- yalnızca çok kısa, atomik bağlamda */
udelay(50);                                 /* 50 mikrosaniye */
mdelay(5);                                  /* 5 milisaniye (kaçın!) */

/* UYUYAN bekleme (tercih edilen, süreç bağlamında) -- CPU'yu bırakır */
msleep(100);                                /* 100 ms uyu */
usleep_range(50, 100);                      /* 50-100 us esnek uyku */
\end{lstlisting}

\uyari{\code{udelay}/\code{mdelay} CPU'yu \emph{meşgul} tutarak bekler (busy-wait)
--- işlemciyi boşa harcar. Sadece atomik bağlamda ve çok kısa süreler için
kullan. Süreç bağlamında her zaman \code{msleep}/\code{usleep\_range} tercih
et; bunlar CPU'yu başka işlere bırakır.}

\gsub{Çekirdek Zamanlayıcıları}

Belirli bir süre sonra bir fonksiyonun çağrılmasını istiyorsak zamanlayıcı
(timer) kurarız --- örneğin bir donanım yanıtı için zaman aşımı.

\begin{lstlisting}
#include <linux/timer.h>

static struct timer_list my_timer;

/* Zaman dolduğunda çağrılır (softirq bağlamı -> uyuyamaz) */
static void my_timer_callback(struct timer_list *t)
{
    pr_info("zamanlayıcı ateşlendi\n");
    /* Periyodik istiyorsak yeniden kur: */
    mod_timer(&my_timer, jiffies + HZ);     /* 1 saniye sonra tekrar */
}

static int __init my_init(void)
{
    timer_setup(&my_timer, my_timer_callback, 0);   /* ilklendir */
    mod_timer(&my_timer, jiffies + HZ);             /* 1 sn sonra ateşle */
    return 0;
}

static void __exit my_exit(void)
{
    del_timer_sync(&my_timer);              /* zamanlayıcıyı güvenle durdur */
}
\end{lstlisting}

\ipucu{Daha yüksek çözünürlük gerekiyorsa \code{hrtimer} (high-resolution
timer) nanosaniye hassasiyeti sunar. Periyodik ve uyuyabilen işler için ise
\code{delayed\_work} (gecikmeli workqueue) genelde en pratik çözümdür ---
hem zamanlar hem de süreç bağlamında çalışır.}

% #####################################################################
\gpart{4}{Aygıt Modeli ve Sürücü Çerçeveleri}
% #####################################################################

% =====================================================================
\gsec{sysfs ile Kullanıcı Arayüzü}
% =====================================================================

sysfs (\code{/sys}), çekirdeğin iç durumunu dosya-benzeri bir hiyerarşi olarak
kullanıcı alanına sunar. Bir sürücü, \emph{nitelik} (attribute) dosyaları
oluşturarak ayarlarını ve durumunu \code{cat}/\code{echo} ile erişilebilir
kılar --- \code{ioctl}'e göre çok daha keşfedilebilir ve betiklenebilir.

\gsub{sysfs Niteliği Oluşturma}

\begin{lstlisting}
/* Kullanıcı dosyayı OKUDUĞUNDA çağrılır (cat) */
static ssize_t value_show(struct device *dev,
                          struct device_attribute *attr, char *buf)
{
    return sysfs_emit(buf, "%d\n", my_value);   /* güvenli tampon yazımı */
}

/* Kullanıcı dosyaya YAZDIĞINDA çağrılır (echo) */
static ssize_t value_store(struct device *dev, struct device_attribute *attr,
                           const char *buf, size_t count)
{
    int ret = kstrtoint(buf, 10, &my_value);    /* metni int'e çevir */
    if (ret < 0)
        return ret;                             /* geçersiz girdi */
    return count;                               /* tüketilen bayt sayısı */
}

/* Nitelik: 0644 izinli, "value" adlı, okuma+yazma dosyası */
static DEVICE_ATTR_RW(value);

/* Kaydet (device_create sonrası): -> /sys/.../mydev/value */
device_create_file(my_device, &dev_attr_value);
\end{lstlisting}

\begin{lstlisting}[style=shstyle]
# Artık kabuktan doğrudan erişilir
$ cat /sys/class/myclass/mydev/value
42
$ echo 100 | sudo tee /sys/class/myclass/mydev/value
\end{lstlisting}

\notk[sysfs kuralı]{sysfs dosyaları ideal olarak \textbf{dosya başına tek bir
değer} içermelidir (bir sayı, bir metin) --- karmaşık ikili yapılar değil.
``Bir dosya, bir değer'' kuralı, sysfs'i basit ve betiklenebilir tutar.
Karmaşık veri için \code{debugfs} veya karakter aygıtı daha uygundur.}

% =====================================================================
\gsec{Platform Sürücüleri}
% =====================================================================

PCI veya USB gibi kendini tanıtan (enumerable) veri yollarında aygıtlar
otomatik keşfedilir. Ama bir SoC üzerindeki tümleşik çevre birimleri
(GPIO, I2C denetleyici, UART\dots) kendini tanıtmaz --- bunlara \emph{platform
aygıtı} denir ve varlıkları sisteme ayrıca bildirilir. Gömülü Linux'ta en çok
karşılaşılan sürücü türüdür.

\gsub{Platform Sürücüsünün İskeleti}

Çekirdek, aygıt ile sürücüyü bir \emph{isim} veya \emph{uyumluluk} dizesiyle
eşleştirir; eşleşince sürücünün \code{probe} fonksiyonu çağrılır.

\begin{lstlisting}
#include <linux/platform_device.h>
#include <linux/mod_devicetable.h>

/* Aygıt bulunup sürücüye bağlandığında çağrılır -- init işini burada yap */
static int my_probe(struct platform_device *pdev)
{
    struct device *dev = &pdev->dev;
    dev_info(dev, "aygıt bulundu, başlatılıyor\n");

    /* Kaynakları al (bellek, IRQ, saat...) -- sonraki bölüm */
    /* devm_* ile ayır: remove'da otomatik temizlenir */
    return 0;
}

/* Aygıt kaldırıldığında çağrılır -- temizlik */
static void my_remove(struct platform_device *pdev)
{
    dev_info(&pdev->dev, "aygıt kaldırıldı\n");
}

/* Device tree "compatible" eşleştirme tablosu */
static const struct of_device_id my_of_match[] = {
    { .compatible = "mycompany,mydevice" },
    { /* sonlandırıcı boş girdi */ }
};
MODULE_DEVICE_TABLE(of, my_of_match);

static struct platform_driver my_driver = {
    .probe  = my_probe,
    .remove = my_remove,
    .driver = {
        .name           = "my_driver",
        .of_match_table = my_of_match,
    },
};

/* Tüm init/exit kalıbını tek makroyla üret */
module_platform_driver(my_driver);
\end{lstlisting}

\bilgi{\code{module\_platform\_driver()} makrosu, \code{module\_init}/
\code{module\_exit} fonksiyonlarını ve sürücü kayıt/kayıt-silme çağrılarını
otomatik üretir --- basit sürücülerde onlarca satır boilerplate'i tek satıra
indirir. Benzer makrolar \code{module\_i2c\_driver}, \code{module\_pci\_driver}
için de vardır.}

% =====================================================================
\gsec{Device Tree ve Kaynaklar}
% =====================================================================

Device Tree (DT), donanımın \emph{nerede} olduğunu (adresler, kesme numaraları,
saatler) koddan ayrı, bildirimsel bir dosyada tarif eder. Böylece aynı sürücü,
farklı kartlarda farklı adreslerle --- kodu değiştirmeden --- çalışır. Gömülü
sistemlerde donanım tanımının standart yoludur.

\gsub{Device Tree Düğümü}

\begin{lstlisting}[style=shstyle]
/* Device Tree kaynağı (.dts) -- donanımı tarif eder, kod değil */
mydevice@10000000 {
    compatible = "mycompany,mydevice";  /* sürücüyle eşleşen dize */
    reg = <0x10000000 0x1000>;          /* register adresi + boyut */
    interrupts = <0 42 4>;              /* kesme numarası ve tipi */
    clocks = <&clk_controller 5>;       /* kullandığı saat */
    my-custom-property = <100>;         /* sürücüye özel ayar */
};
\end{lstlisting}

\gsub{probe İçinde Kaynak Alma}

Sürücü, DT'de tanımlı kaynakları \code{probe} sırasında okur. \code{devm\_*}
sürümleri, aygıt kaldırılınca kaynakları otomatik serbest bırakır.

\begin{lstlisting}
static int my_probe(struct platform_device *pdev)
{
    struct device *dev = &pdev->dev;
    void __iomem *regs;
    int irq, ret;
    u32 custom;

    /* 1) Register bellek bölgesini al ve eşle (MMIO) */
    regs = devm_platform_ioremap_resource(pdev, 0);
    if (IS_ERR(regs))
        return PTR_ERR(regs);

    /* 2) Kesme numarasını al */
    irq = platform_get_irq(pdev, 0);
    if (irq < 0)
        return irq;

    ret = devm_request_irq(dev, irq, my_irq_handler, 0, "mydev", pdev);
    if (ret)
        return ret;

    /* 3) DT'den özel bir özelliği oku */
    if (of_property_read_u32(dev->of_node, "my-custom-property", &custom) == 0)
        dev_info(dev, "özel değer: %u\n", custom);

    return 0;
}
\end{lstlisting}

\gsub{Donanım Register'larına Erişim (MMIO)}

Çoğu çevre birimi, register'ları bellek adreslerine eşlenmiş olarak sunar
(Memory-Mapped I/O). Bu register'lara \emph{asla} doğrudan işaretçiyle değil,
özel erişim fonksiyonlarıyla (\code{readl}/\code{writel}) erişilir --- bunlar
doğru bariyerleri ve sıralamayı garanti eder.

\begin{lstlisting}
#include <linux/io.h>

#define REG_CONTROL  0x00                   /* register ofsetleri */
#define REG_STATUS   0x04
#define REG_DATA     0x08

/* 32-bit register oku */
u32 status = readl(regs + REG_STATUS);

/* 32-bit register'a yaz */
writel(0x1, regs + REG_CONTROL);            /* aygıtı etkinleştir */

/* Bir durum bitini bekle (zaman aşımıyla) */
u32 val;
ret = readl_poll_timeout(regs + REG_STATUS, val,
                         (val & BIT(0)),    /* koşul: bit 0 set olana kadar */
                         10, 1000);         /* 10us aralık, 1000us zaman aşımı */
if (ret)
    dev_err(dev, "aygıt hazır olmadı\n");
\end{lstlisting}

\uyari{Donanım register'larına normal işaretçiyle (\code{*(u32*)addr})
erişmek yanlıştır: derleyici erişimleri yeniden sıralayabilir, önbelleğe
alabilir veya birleştirebilir --- donanımla iletişimi bozar. Her zaman
\code{readl}/\code{writel} (ve 8/16 bit için \code{readb}/\code{readw})
kullan; bunlar \code{volatile} erişim ve bellek bariyerlerini garanti eder.}

% #####################################################################
\gpart{5}{İleri Konular ve Hata Ayıklama}
% #####################################################################

% =====================================================================
\gsec{Bekleme Kuyrukları ve Bloklamalı G/Ç}
% =====================================================================

Çoğu zaman \code{read}, veri henüz hazır değilken çağrılır. Sürücü ya hemen
``veri yok'' dönmeli (bloklamasız) ya da veri gelene kadar süreci
\emph{uyutmalıdır} (bloklamalı). İkincisi için \emph{bekleme kuyrukları}
(wait queues) kullanılır: süreç uyur, CPU'yu harcamaz; veri gelince uyandırılır.

\gsub{Bekleme Kuyruğu ile Uyutma ve Uyandırma}

\begin{lstlisting}
#include <linux/wait.h>
#include <linux/sched.h>

static DECLARE_WAIT_QUEUE_HEAD(read_queue);
static bool data_ready;

/* OKUMA tarafı: veri yoksa uyu */
static ssize_t my_read(struct file *file, char __user *buf,
                       size_t count, loff_t *offset)
{
    /* Bloklamasız istendiyse ve veri yoksa hemen dön */
    if ((file->f_flags & O_NONBLOCK) && !data_ready)
        return -EAGAIN;

    /* data_ready true olana kadar UYU (sinyalle kesilebilir) */
    if (wait_event_interruptible(read_queue, data_ready))
        return -ERESTARTSYS;                /* sinyal geldi */

    /* Buraya gelindi: veri hazır -> kullanıcıya kopyala */
    if (copy_to_user(buf, device_buffer, count))
        return -EFAULT;
    data_ready = false;
    return count;
}

/* ÜRETİCİ taraf (ör. kesme işleyici): veri hazır, uyuyanları uyandır */
static void data_arrived(void)
{
    data_ready = true;
    wake_up_interruptible(&read_queue);     /* bekleyen okuyucuları uyandır */
}
\end{lstlisting}

\bilgi{\code{wait\_event\_interruptible} bir \emph{koşul} bekler ve koşul
sağlanana dek uyur; sahte uyanmalara (spurious wakeup) karşı koşulu döngüde
tekrar kontrol eder --- bu yüzden onu elle döngüye almana gerek yoktur.
``interruptible'' sürümü, sürecin sinyalle (Ctrl+C) uyandırılabilmesini sağlar;
bu neredeyse her zaman tercih edilir (aksi halde süreç öldürülemez).}

\gsub{poll/select Desteği}

Bir uygulamanın birden çok aygıtı aynı anda \code{poll}/\code{select}/
\code{epoll} ile beklemesi için sürücü \code{poll} fonksiyonunu sağlamalıdır.

\begin{lstlisting}
#include <linux/poll.h>

static __poll_t my_poll(struct file *file, struct poll_table_struct *wait)
{
    __poll_t mask = 0;

    poll_wait(file, &read_queue, wait);     /* kuyruğu poll tablosuna kaydet */

    if (data_ready)
        mask |= POLLIN | POLLRDNORM;        /* okunacak veri var */

    return mask;                            /* hazır olan olayları bildir */
}
/* file_operations: .poll = my_poll, */
\end{lstlisting}

% =====================================================================
\gsec{Hata Ayıklama Teknikleri}
% =====================================================================

Çekirdekte hata ayıklamak kullanıcı alanından zordur: \code{gdb} ile adım adım
ilerlemek genelde pratik değildir. Bu yüzden çekirdek, kendine özgü zengin
hata ayıklama araçları sunar.

\begin{center}
\begin{tabularx}{\textwidth}{@{}>{\bfseries\color{anaMavi}}l X@{}}
\toprule
\rowcolor{acikMavi}\color{anaMavi}\textbf{Araç} & \color{anaMavi}\textbf{Kullanım}\\
\midrule
dmesg / printk & En temel; mesajlarla izleme.\\
\rowcolor{acikGri} dynamic debug & \code{pr\_debug}'ı çalışırken aç/kapa.\\
ftrace & Fonksiyon çağrı izleme, zamanlama.\\
\rowcolor{acikGri} kgdb / kdb & Çekirdek düzeyinde kesme noktalı hata ayıklama.\\
KASAN & Bellek hatası dedektörü (use-after-free, taşma).\\
\rowcolor{acikGri} debugfs & Hata ayıklama için serbest biçimli sysfs.\\
\bottomrule
\end{tabularx}
\end{center}

\gsub{Kernel Oops ve Panic Okuma}

Bir sürücü hatası (örneğin \code{NULL} işaretçi dereference), \emph{oops}
üretir --- çekirdek bir yığın izi (stack trace) ve register durumu basar.
Bunu okumak, hatanın nerede olduğunu bulmanın anahtarıdır.

\begin{lstlisting}[style=shstyle]
# Tipik bir oops mesajı (dmesg'de)
BUG: kernel NULL pointer dereference at 0000000000000000
RIP: 0010:my_read+0x1a/0x80 [mymod]     # <- HATA BURADA: my_read, +0x1a ofset
Call Trace:
 vfs_read+0x9d/0x150                     # çağrı yığını (aşağıdan yukarı okunur)
 ksys_read+0x67/0xe0
 do_syscall_64+0x59/0x90

# Ofseti kaynak satırına çevir (hata ayıklama sembolleriyle)
$ addr2line -e mymod.ko 0x1a
/home/user/mymod.c:42                    # tam olarak 42. satır!
\end{lstlisting}

\gsub{debugfs ile İç Durum Sunma}

\begin{lstlisting}
#include <linux/debugfs.h>

static struct dentry *debug_dir;
static u32 error_count;                     /* izlemek istediğimiz sayaç */

static int __init my_init(void)
{
    /* /sys/kernel/debug/mydev/ altında dosyalar oluştur */
    debug_dir = debugfs_create_dir("mydev", NULL);
    debugfs_create_u32("error_count", 0644, debug_dir, &error_count);
    return 0;
}

static void __exit my_exit(void)
{
    debugfs_remove_recursive(debug_dir);    /* tümünü temizle */
}
\end{lstlisting}

\ipucu{Geliştirme sırasında \code{debugfs} paha biçilmezdir: bir sayaç veya
bayrağı, tek satırla (\code{debugfs\_create\_u32}) kabuktan izlenebilir hale
getirir. sysfs'in aksine debugfs'in ``kararlı ABI'' garantisi yoktur --- yani
kuralları gevşektir ve tam da hata ayıklama için tasarlanmıştır. Üretim
arayüzleri için sysfs, hata ayıklama için debugfs kullan.}

% =====================================================================
\gsec{Kapanış: Sürücü Geliştirme Yol Haritası}
% =====================================================================

Bu rehberde bir Linux sürücüsünün tüm temel parçalarını gördük. Gerçek bir
sürücü, bu parçaların belirli bir donanım için birleştirilmesidir.

\begin{center}
\begin{tabularx}{\textwidth}{@{}>{\bfseries\color{anaMavi}}l X@{}}
\toprule
\rowcolor{acikMavi}\color{anaMavi}\textbf{İhtiyaç} & \color{anaMavi}\textbf{Kullanılan mekanizma}\\
\midrule
Modülü yükle/kaldır & \code{module\_init} / \code{module\_exit}\\
\rowcolor{acikGri} Kullanıcı alanına arayüz & Karakter aygıtı (\code{cdev} + \code{file\_operations})\\
Komut gönderme & \code{ioctl} veya sysfs niteliği\\
\rowcolor{acikGri} Otomatik \code{/dev} düğümü & \code{class\_create} + \code{device\_create}\\
Bellek ayırma & \code{kzalloc} / \code{devm\_kzalloc}\\
\rowcolor{acikGri} Paylaşılan veri koruma & \code{mutex} / \code{spinlock} / \code{atomic\_t}\\
Donanım olayları & Kesme işleyici + alt yarı (workqueue)\\
\rowcolor{acikGri} Zaman aşımı / periyodik iş & Zamanlayıcı / \code{delayed\_work}\\
Gömülü donanım keşfi & Platform sürücüsü + Device Tree\\
\rowcolor{acikGri} Bloklamalı okuma & Bekleme kuyruğu + \code{poll}\\
Hata ayıklama & \code{dmesg}, \code{ftrace}, \code{debugfs}, KASAN\\
\bottomrule
\end{tabularx}
\end{center}

\ipucu{Sürücü geliştirmeyi öğrenmenin en iyi yolu \textbf{gerçek kod okumaktır}:
çekirdek kaynağındaki \code{drivers/} dizini, her tür donanım için binlerce
örnek sürücü içerir. Basit bir örnekle (\code{drivers/char/} altında) başla,
onu bir sanal makinede derleyip değiştir. Ayrıca çekirdek belgeleri
(\code{Documentation/} dizini) ve ``Linux Device Drivers'' (LDD3) kitabı
temel başvurulardır. Küçük bir karakter sürücüsü yaz, \code{dmesg} ile izle,
kır ve düzelt --- öğrenmenin tek yolu budur. İyi kodlamalar!}

\vfill
\begin{center}
{\color{ortaGri}\rule{0.5\linewidth}{0.6pt}}\\[6pt]
{\color{ortaGri}\small Rehberin sonu --- Çekirdek alanında dikkatli olun!}
\end{center}

\end{document}
